\scanpage{14}
\uuline{Non}, car ce groupe est le $=$ pour \deleted{la str}
$\omega=\{u,u^{-1}\}$ et $\omega^{(r)}=\{u^r,u^{-r}\}$ où $r\in\mathbb Z$, $(r,n)=1$.

Il faut donc en plus se donner une
partie
\[
\omega\subset\Gamma\tag*{(4)}
\]
avec
\[
\left\{\begin{aligned}
&\operatorname{card}(\omega)=2,\qquad\omega=\{u,u^{-1}\}\\
&\deleted{\unsure{\Gamma|\Gamma^0}}
\end{aligned}\right.\tag*{(5)}
\]
\deleted{\unsure{et|de}} telle que, si $\Gamma^{\unsure{0|\circ}}$ est le ss-gpe de $\Gamma$
engendré par $\omega$ (i.e. par \deleted{l'}un \uline{\unsure{seul|quelq.}} \deleted{\unsure{élémt}}
de ses éléments) on ait
\[
\left\{\begin{aligned}
&\operatorname{Card}\Gamma/\Gamma^0=2\\
&\Gamma/\Gamma^0\ \text{opère sur }\Gamma^0\ \text{par}\quad\lambda\longmapsto\lambda^{-1}\\
&\text{l'ext. }\Gamma\ \text{de }\Gamma/\Gamma^0\ \text{par }\Gamma^0\ \text{est semi-d.te,}\\
&\hspace{4em}\text{i.e. }\exists\,\sigma\in\Gamma,\ \sigma^2=1,\ \sigma\longmapsto\text{él\textsuperscript{t}}
\\&\hfill {-1}\ \text{de }\Gamma/\unsure{\Gamma_0|\Gamma_1|\Gamma^0}
\end{aligned}\right.\tag*{(6)}
\]
[\uline{NB} Ces conditions, qui précisent (1) \inserted{(modulo la condition que $\Gamma^0$ est fini)}
doivent bien sûr remplacer (2)]

En résumé, $P$ est déterminé
par la donnée

a) du ss-gpe $\Gamma\subset\mathit{Aff}(E)$

b) de la partie $\omega\subset\Gamma$
satisfaisant (5) (6)

c) de $S\in\Gamma\backslash E$ (ens des orbites)
% archive sheet number: 13

